\abstractPT  % Do NOT modify this line

Play4Change é uma plataforma de aprendizagem adaptativa que combina inteligência artificial generativa, mecânicas de gamificação e revisão por pares para transformar conteúdo educativo em desafios diários personalizados. A plataforma responde a duas necessidades críticas: a educação para a sustentabilidade alinhada com os Objetivos de Desenvolvimento Sustentável das Nações Unidas, e a literacia digital.

Do ponto de vista técnico, o sistema é constituído por uma aplicação móvel multiplataforma em Kotlin Multiplatform (KMP) com Decompose e o padrão MVI para iOS e Android; um servidor Spring Boot em Kotlin com Design Orientado a Domínio, arquitetura hexagonal e Clean Architecture; um módulo Kotlin partilhado que garante contratos de API type-safe entre cliente e servidor; e uma página de destino em TypeScript/React com painel de administração. A geração de conteúdo é realizada por um pipeline de IA que combina o modelo Mistral (via LangChain4j) com Retrieval-Augmented Generation (RAG) e pesquisa por similaridade vetorial com pgvector. Todas as decisões arquiteturais estão documentadas em 16 Architecture Decision Records.

O sistema privilegia custos sustentáveis e previsíveis, autenticação sem palavra-passe (magic link + OAuth 2.0), cache Redis com degradação silenciosa, observabilidade completa com Micrometer e Grafana, e preparação para Kubernetes com gatilhos de migração explicitamente definidos. A revisão por pares das tarefas aplicadas oferece avaliação sem custo marginal de IA, gerando simultaneamente eventos de aprendizagem para os revisores.

\keywords{Aprendizagem Adaptativa \sep Gamificação \sep Inteligência Artificial Generativa \sep Kotlin Multiplatform \sep Sustentabilidade \sep Clean Architecture \sep Arrow Either \sep pgvector}
